\documentclass[11pt]{article}
\usepackage[T1]{fontenc}
\usepackage{amsmath,amssymb,amsthm,geometry,hyperref}
\geometry{margin=1in}
\newtheorem{proposition}{Proposition}
\newtheorem{lemma}{Lemma}
\newcommand{\Tr}{\operatorname{Tr}}
\title{AIM 001: cylinder reduction and the missing geometric estimate}
\author{Autonomous research programme: senior run 1}
\date{8 September 2026}
\begin{document}
\maketitle

\paragraph{Status.} This is a cumulative starting note, not a candidate solution.
The exact target is both P\'olya inequalities for every bounded connected
Lipschitz domain in every dimension $d\geq2$, at every eigenvalue index.
The results below are exact reductions and method obstructions; no new
inequality for arbitrary actual domains is claimed, and novelty is not claimed.
Source-scope checks and detailed countermodel proofs are separate accompanying
artifacts. Earlier off-repository research has not been imported as evidence.

\section{Normalization and endpoints}
Let $C_d=\omega_d/(2\pi)^d$ and $V=|\Omega|$. Dirichlet eigenvalues are
$\lambda_1,\lambda_2,\ldots$, while Neumann eigenvalues are
$\mu_0=0,\mu_1,\ldots$, always with multiplicity. Put
\[
 N_D(E)=\#\{j\geq1:\lambda_j\leq E\},\qquad
 N_N(E)=\#\{j\geq0:\mu_j<E\}.
\]
The target is equivalently $N_D(E)\leq C_dVE^{d/2}\leq N_N(E)$ for all
$E\geq0$. For example, writing $m=N_N(E)$ for $E>0$, one has
$\mu_m\geq E$ and thus $C_dVE^{d/2}\leq m$ from the eigenvalue formulation.
Conversely a failure $\mu_k>(k/(C_dV))^{2/d}$ allows an energy strictly
between these values, contradicting the counting formulation.
The Dirichlet equivalence follows by evaluating at eigenvalues, including
the last index in each multiplicity block. At $E=0$ both inequalities are
trivial. Strict/inclusive counting must not be interchanged at a fixed energy.

\section{Actual-domain cylinder bridge}
For $E\geq0$ set $R_{1/2}^D(E)=\sum_j(E-\lambda_j)_+^{1/2}$ and define
$R_{1/2}^N$ using all $\mu_j$, including $\mu_0$. All these sums are finite
for a bounded Lipschitz domain. Write $I_L=(0,L)$.

\begin{proposition}
For a fixed bounded connected Lipschitz $\Omega\subset\mathbb R^d$, the
Dirichlet P\'olya inequality on $\Omega\times I_L$ for every $L>0$ is
equivalent to
\[
 R_{1/2}^D(E)\leq\pi C_{d+1}VE^{(d+1)/2}\quad(E\geq0).
\]
The corresponding Neumann assertion is equivalent to the reversed
half-Riesz inequality, with the zero mode included.
\end{proposition}

\begin{proof}
The Dirichlet and Neumann quadratic forms on the product give the tensor
sum spectra $\lambda_j+\pi^2n^2/L^2$, $n\geq1$, and
$\mu_j+\pi^2n^2/L^2$, $n\geq0$, respectively. This follows by taking the
tensor products of the complete form eigenbases; smooth boundary or $H^2$
regularity is not required. Products of bounded Lipschitz domains with an
interval remain bounded connected Lipschitz domains. Consequently
\[
 N_{D,L}(E)=\sum_{\lambda_j\leq E}
 \left\lfloor\frac L\pi\sqrt{E-\lambda_j}\right\rfloor,
\quad
 N_{N,L}(E)=\sum_{\mu_j<E}
 \left\lceil\frac L\pi\sqrt{E-\mu_j}\right\rceil.
\]
The Neumann formula counts nonnegative integers strictly less than a
positive real number, which explains both the ceiling and the restriction.
Thus $0\leq LR_{1/2}^D/\pi-N_{D,L}\leq N_D(E)$ and
$0\leq N_{N,L}-LR_{1/2}^N/\pi\leq N_N(E)$.
The signs imply sufficiency. For necessity divide by $L$ and let $L$ tend
to infinity at each fixed $E$: the finite rounding bounds prove the limit.
There is no exchange with an unbounded energy or spectral-index limit.
Finally the volume of the cylinder is $LV$ and
\[
 \pi C_{d+1}=\frac{\Gamma(3/2)}{(4\pi)^{d/2}\Gamma((d+3)/2)},
\]
the sharp half-Riesz coefficient.
\end{proof}

\paragraph{Explicit finite transfer.}
If a Dirichlet excess $\delta>0$ at energy $E$ is witnessed using only
$m$ base eigenvalues, that is
\[
 \sum_{j=1}^m(E-\lambda_j)_+^{1/2}
 >\pi C_{d+1}VE^{(d+1)/2}+\delta,
\]
then every $L>\pi m/\delta$ yields a cylinder counterexample by
$\lfloor x\rfloor\geq x-1$. A Neumann deficit $\delta>0$ transfers for
$L>\pi N_N(E)/\delta$, using $\lceil x\rceil\leq x+1$.
For a computational Dirichlet certificate, certified upper bounds on the
chosen eigenvalues give lower bounds on the sum; an upper bound on volume
is safe. For a Neumann certificate, an upper bound on the entire half-Riesz
sum and a certified tail cutoff are necessary; lower bounds on finitely
many eigenvalues alone do not control omitted modes. Energy powers, square
roots, and constants must have exact or outward-rounded enclosures.
No such actual-domain violation or certificate has been obtained.

In a planar base the coefficient is $V/(6\pi)$. The next concrete
intermediate target is therefore $R_{1/2}^D(E)\leq VE^{3/2}/(6\pi)$,
or its Neumann reverse. It concerns actual domains and would settle a
nontrivial subclass of the three-dimensional target. It would not settle
arbitrary three-dimensional domains or the full repository problem.

\section{A geometric identity and an impossible stronger inequality}
Let $U:L^2(\Omega)\to L^2(\mathbb R^d)$ be zero extension and use the
unitary Fourier transform $\mathcal F$. For $E>0$ define
\[
 A_E=1_{\{|\xi|^2\leq E\}}\mathcal F U,\qquad T_E=A_E^*A_E.
\]
The integral kernel of $A_E$ has constant squared modulus $(2\pi)^{-d}$,
so it is Hilbert--Schmidt, $0\leq T_E\leq I$, and
$\Tr T_E=C_dVE^{d/2}$. Let $P_E$ be the finite-rank Dirichlet projection
onto eigenvalues at most $E$. Trace cyclicity for trace-class operators
gives the exact identity
\[
 C_dVE^{d/2}-N_D(E)
 =\Tr((I-P_E)T_E)-\Tr(P_E(I-T_E)).                 \tag{1}
\]
Thus the desired counting inequality compares the Fourier mass inside the
ball contributed by high Dirichlet modes against the Fourier mass outside
the ball contributed by low modes. Both traces are well defined and
nonnegative. Parseval and Tonelli identify the first trace with the
convergent sum of high-mode masses inside the ball.

\begin{lemma}
If $P_E\ne0$, then $T_E\geq P_E$ is impossible.
\end{lemma}
\begin{proof}
For nonzero $f\in L^2(\Omega)$, zero extension is integrable and compactly
supported. Its Fourier transform is real analytic. If
$\langle f,(I-T_E)f\rangle=0$, that transform vanishes almost everywhere
outside the ball; continuity and analytic uniqueness imply it vanishes
identically, a contradiction. Hence this quadratic form is strictly
positive on every nonzero $f$. In particular it is positive for a nonzero
$f$ in the range of $P_E$, contradicting $T_E\geq P_E$.
\end{proof}

Equation (1) is a reformulation, not a bound or a claimed new result about
domains. It shows why a direct operator domination argument fails and
identifies the compensation that a successful geometric proof must control.
Zero-extended Dirichlet eigenfunctions belong to $H^1(\mathbb R^d)$, so
their Fourier second moments equal their energies; this weighted control
does not by itself give the unweighted trace comparison (1). No unproved
boundary-normal regularity or interchange of infinite energy sums is used.

\section{Adversarial results and remaining scope}
The accompanying exact notes establish the following logical obstructions.
\begin{itemize}
\item The sequence $a_j=j+\lceil\sqrt j\rceil$ with entries $90$ through
$100$ lowered to $399/4$ violates its planar pointwise Weyl bound, despite
sharp heat and sharp Riesz bounds of every order at least $1/128$, and the counting asymptotic
$N(E)=E-\sqrt E+O(1)$. The accompanying referee note quantifies how far
below order one the same example remains a valid obstruction.
\item A different mean-flattened block of $j+1$ preserves every sharp
Riesz bound of order at least one and the full planar Yang inequality,
but violates the sharp half-Riesz bound by an explicitly certified margin.
Its proof and rational verifier are in the spectral-attack artifacts.
An independently audited asymptotic extension starts from square eigenvalues
and preserves the exact eventual square spectrum as well as Yang and these
convex trace bounds. It uses the stated external positive-order two-term
Riesz theorem and supplies no effective block threshold.
\end{itemize}
These are abstract sequences, not Laplacian spectra. They exclude
inferences based only on the stated scalar inequalities and asymptotics.
They do not exclude a proof using additional geometric information, and
they are not counterexamples to P\'olya. Bounds for one fixed positive
Riesz order cannot simply be differentiated or sent to order zero. If a
single spectrum satisfied all positive orders down to zero, the fixed
energy limit would recover counting; the cutoff-dependent countermodels
do not contradict this fact.

\paragraph{Source baseline.}
The source audit records the precise scope of balls, planar annuli,
convex-domain positive-order improvements, and epsilon-loss high-energy
bounds. In particular, the explicit threshold in
\href{https://arxiv.org/abs/2507.04307}{Jiang--Lin} depends on epsilon;
it gives no exact all-index reduction by setting epsilon to zero.
\href{https://arxiv.org/abs/2407.11808}{Frank--Larson's two-term Riesz
asymptotics} apply to positive order and are not automatically counting
asymptotics. The standard product connection may also be compared with
\href{https://arxiv.org/abs/2402.12093}{He--Wang's thin-products work}.

\section{Concurrent Sol audit and the updated first bottleneck}
Sol attempt 1 committed during this run and is preserved unchanged. Its
full $k=1,2$ chain follows from the four established shape inequalities;
the source audit checks Bucur--Henrot's indexing and Lipschitz scope.
Its criterion $k\leq r\beta_d$, where
$\beta_d=j_{d/2-1,1}^d\omega_d^2/(2\pi)^d$ and $r$ is the actual nodal
count, is valid. The universal band $2\leq k\leq\lfloor2\beta_d\rfloor$
also follows directly from the second-eigenvalue bound and
$\lambda_k\geq\lambda_2$, with no nodal-boundary regularity assumption.
The independent polynomial verifier confirms the unit-disk trial quotients
$6,16,16$ and the strict failure of the stronger equal-three-ball proposal.

The generic stability route has an existing external theorem:
\href{https://arxiv.org/html/2506.05870v3}{de Villeroch\'e, Theorem 1.1}.
Normalize area to $\pi$, set $A=2j_{0,1}^2<12$ and
$B=\lambda_3(\Theta)=2\lambda_2(D)>12$, with $\Theta$ two equal disks.
That theorem implies that any $\lambda_3(\Omega)\leq12$ must have
\[
 \lambda_2(\Omega)-A\geq
 \left(\frac{B-12}{81C_2\,12^{17/18}}\right)^{18}=:\eta>0.
\]
The separate stability addendum proves this consequence and records its
external dependency. No usable numerical $C_2$ was obtained; no comparison
of $\eta$ with the whole width $12-A$ is asserted. The stronger
positive-part exponent in that paper has the wrong direction for this
lower-bound task.

\paragraph{Campaign decision.}
Continue problem 001 in research phase, now incorporating Sol attempt 1.
Prioritize quantitative lower-side planar third-mode stability, followed
by the Neumann third mode; retain half-Riesz as a secondary route.
There is no complete proof
candidate and no publishable actual-domain partial result established by
this run. Subsequent work must quantify the surviving third-mode interval
or supply a new geometric premise for the
half-Riesz inequality or (1), or a rigorous actual-domain violation; more
scalar desmoothing or larger unvalidated numerical tables do not address
the identified gap. The full Dirichlet and Neumann target remains open.
\end{document}
